\documentclass[11pt]{article}
\usepackage[top=1.00in, bottom=1.0in, left=1in, right=1in]{geometry}
\renewcommand{\baselinestretch}{1.1}
\usepackage{graphicx}
\usepackage{natbib}
\usepackage{amsmath}

\def\labelitemi{--}
\parindent=0pt

\begin{document}
\bibliographystyle{/Users/Lizzie/Documents/EndnoteRelated/Bibtex/styles/besjournals}
\renewcommand{\refname}{\CHead{}}

\setlength{\parindent}{0cm}
\setlength{\parskip}{5pt}

% just a cheap file to fix math errors more easily...

\begin{align}
  \label{modely}
  y_{i,j} \sim \mathcal{N}(\mu_j, \sigma_e^2)
\end{align}
where
\begin{align}
  \label{modelmu}
  \mu_j = \alpha_j 
  \end{align}
and $\sigma_e^2$ represents random error unrelated to the phylogeny. And the species-specific intercept $\alpha_j$ is elements of the following normal random vectors:

\begin{align}
    \label{phybetas}
  \boldsymbol{\alpha} = [\alpha_1, \ldots, \alpha_n]^T & \text{ such that }
  \boldsymbol{\alpha} \sim \machtal{N}(\mu_{\alpha},\boldsymbol{\Sigma_{\alpha}}) 
\end{align}


\begin{align}
  \label{phymatnew}
\begin{bmatrix}
  \sigma^2 & \lambda \times \sigma^2 \times \rho_{12} & \ldots & \lambda \times \sigma^2 \times \rho_{1n} \\
  \lambda \times \sigma^2 \times \rho_{21} & \sigma^2 & \ldots & \lambda \times \sigma^2 \times \rho_{2n} \\
  \vdots & \vdots & \ddots & \vdots \\
  \lambda \times \sigma^2 \times \rho_{n1} & \lambda \times \sigma^2 \times \rho_{n2} & \ldots & \sigma^2\\
\end{bmatrix}
\end{align}

\begin{align}
  \label{phymatalt}
\begin{bmatrix}
  \sigma^2 & \lambda \times \sigma_{12} & \ldots & \lambda \times \sigma_{1n} \\
  \lambda \times \sigma_{21} & \sigma^2 & \ldots & \lambda \times \sigma_{2n} \\
  \vdots & \vdots & \ddots & \vdots \\
  \lambda \times \sigma_{n1} & \lambda \times \sigma_{n2} & \ldots & \sigma^2_i \\
\end{bmatrix}
\end{align}

Above I have dropped $\rho$. Now $\sigma^2$ is still the Brownian motion rate of evolution across the tree, and $\sigma$ is the covariance between representative taxa. This is not an uncommon representation but I think we could have been clearer if we presented the last part of our notation as:


  \begin{equation} {\bf \Sigma}  =  \sigma^2_{BM} \cdot 
\begin{bmatrix}
  \rho_{1,1} & \lambda \cdots \rho_{1,2} & \ldots & \lambda \cdots \rho_{1,n} \\
  \lambda \cdots \rho_{2,1} &  \rho_{2,2} & \ldots & \lambda \cdots \rho_{2,n} \\
  \vdots & \vdots & \ddots & \vdots \\
  \lambda \cdots \rho_{n,1} & \lambda \cdots \rho_{n,2} & \ldots & \rho_{n,n}\\
\end{bmatrix}



\end{document}